% Анотація — за Додатком В навчального посібника (не менше 650 символів).
% Кількості сторінок, рисунків, таблиць, джерел і додатків — з build/main.kpistats попередньої
% збірки (клас kpi-thesis записує їх наприкінці документа; latexmk перезбирає, доки не зійдеться).
\kpiloadstats{main}{build/main.kpistats}
\structural{АНОТАЦІЯ}

Хребтань Максим Віталійович. Курсова робота на тему «\kpiTitle».

Текстова частина курсової роботи складається із вступу, пʼяти розділів, висновків, списку
використаних джерел з \cwStat{sources}{найменування}{найменувань}{найменувань} та
\cwStat{appendices}{додатка}{додатків}{додатків}, які викладено на
\cwStat{pages}{сторінці}{сторінках}{сторінках}. В тексті роботи оформлено
\cwStat{figures}{рисунок}{рисунки}{рисунків}, \cwStat{tables}{таблицю}{таблиці}{таблиць}.

У роботі поставлено завдання дослідити, як пʼять стратегій пріоритетного упорядкування черги за
різних інтенсивностей надходження впливають на середній час перебування завдання в системі,
середню кількість завдань у черзі та середній час очікування в інформаційній системі, що обробляє
завдання трьох типів зі спільною оперативною памʼяттю: завдання A і~B можуть виконуватися
одночасно, а завдання~C монополізує систему.

Імітаційна модель розроблена на основі формального опису системи стохастичною мережею Петрі, у
якій обмеження на одночасне виконання задано кратністю дуг, а стратегії -- пріоритетами та
ймовірностями запуску переходів. Алгоритм імітації реалізовано мовою Java на основі ядра
бібліотеки PetriObjModelPaint, розширеного класами-нащадками; коректність моделі підтверджено
верифікацією. Статистичну обробку виконано мовою Python.

Проведено дослідження: визначено перехідний період, час моделювання й кількість прогонів,
виконано повний факторний експеримент «стратегія $\times$ інтенсивність» (2100~прогонів),
дисперсійний і регресійний аналіз.

Результати моделювання свідчать, що за інтенсивностей з умови задачі система перевантажена за
будь-якої стратегії, а в області стійкості найменші значення всіх трьох показників дають
стратегії з найнижчим пріоритетом завдань~C -- середній час перебування в них до 2,3~раза менший,
ніж у стратегіях з пріоритетом~C, коштом довшого очікування завдань~C.

Зроблено висновки щодо вибору стратегії упорядкування і запропоновано збільшити памʼять системи
до трьох умовних одиниць; пробний запуск показав, що система стає стійкою за інтенсивностей з
умови задачі, а за зменшеного потоку середній час перебування скорочується на 70--86~\%, тож для
потоку, близького до заданого, ця зміна ресурсу доцільна.

\noindent Ключові слова: імітаційна модель, алгоритм імітації, стохастична мережа Петрі,
стратегія упорядкування черги, пріоритетне обслуговування, факторний експеримент, дисперсійний
аналіз, Java.
